\documentclass[10pt,a4paper]{amsart}
\usepackage{amsmath,amssymb,amsthm,graphicx}
\usepackage{verbatim}
\newenvironment{reference}{\comment}{\endcomment}
%\newenvironment{reference}{}{}
\newenvironment{slogan}{\comment}{\endcomment}
\newenvironment{history}{\comment}{\endcomment}

% For commutative diagrams we use Xy-pic
\usepackage[all]{xy}

% We use 2cell for 2-commutative diagrams.
\xyoption{2cell}
% Theorem environments.
%
\theoremstyle{plain}
\newtheorem{theorem}[subsection]{Theorem}
\newtheorem{proposition}[subsection]{Proposition}
\newtheorem{lemma}[subsection]{Lemma}

\theoremstyle{definition}
\newtheorem{definition}[subsection]{Definition}
\newtheorem{example}[subsection]{Example}
\newtheorem{exercise}[subsection]{Exercise}
\newtheorem{situation}[subsection]{Situation}

\theoremstyle{remark}
\newtheorem{remark}[subsection]{Remark}
\newtheorem{remarks}[subsection]{Remarks}

\numberwithin{equation}{subsection}

% Macros
%
\def\lim{\mathop{\mathrm{lim}}\nolimits}
\def\colim{\mathop{\mathrm{colim}}\nolimits}
\def\Spec{\mathop{\mathrm{Spec}}}
\def\Hom{\mathop{\mathrm{Hom}}\nolimits}
\def\Ext{\mathop{\mathrm{Ext}}\nolimits}
\def\SheafHom{\mathop{\mathcal{H}\!\mathit{om}}\nolimits}
\def\SheafExt{\mathop{\mathcal{E}\!\mathit{xt}}\nolimits}
\def\Sch{\mathit{Sch}}
\def\Mor{\mathop{\mathrm{Mor}}\nolimits}
\def\Ob{\mathop{\mathrm{Ob}}\nolimits}
\def\Sh{\mathop{\mathit{Sh}}\nolimits}
\def\NL{\mathop{N\!L}\nolimits}
\def\CH{\mathop{\mathrm{CH}}\nolimits}
\def\proetale{{pro\text{-}\acute{e}tale}}
\def\etale{{\acute{e}tale}}
\def\QCoh{\mathit{QCoh}}
\def\Ker{\mathop{\mathrm{Ker}}}
\def\Im{\mathop{\mathrm{Im}}}
\def\Coker{\mathop{\mathrm{Coker}}}
\def\Coim{\mathop{\mathrm{Coim}}}

% Boxtimes
%
\DeclareMathSymbol{\boxtimes}{\mathbin}{AMSa}{"02}

%
% Macros for moduli stacks/spaces
%
\def\QCohstack{\mathcal{QC}\!\mathit{oh}}
\def\Cohstack{\mathcal{C}\!\mathit{oh}}
\def\Spacesstack{\mathcal{S}\!\mathit{paces}}
\def\Quotfunctor{\mathrm{Quot}}
\def\Hilbfunctor{\mathrm{Hilb}}
\def\Curvesstack{\mathcal{C}\!\mathit{urves}}
\def\Polarizedstack{\mathcal{P}\!\mathit{olarized}}
\def\Complexesstack{\mathcal{C}\!\mathit{omplexes}}
% \Pic is the operator that assigns to X its picard group, usage \Pic(X)
% \Picardstack_{X/B} denotes the Picard stack of X over B
% \Picardfunctor_{X/B} denotes the Picard functor of X over B
\def\Pic{\mathop{\mathrm{Pic}}\nolimits}
\def\Picardstack{\mathcal{P}\!\mathit{ic}}
\def\Picardfunctor{\mathrm{Pic}}
\def\Deformationcategory{\mathcal{D}\!\mathit{ef}}
\usepackage{cleveref}
\crefname{originallemma}{original lemma}{original lemmas}
\Crefname{originallemma}{Original lemma}{Original lemmas}
\makeatletter
\newlabel{lemma-NAK}{{00DV}{source}{}{sourcepointer-1}{}}
\newlabel{lemma-NAK@cref}{{[originallemma][0][]00DV}{[1][0][]source}{}{}{}}
\makeatother
\crefname{originallemma}{original lemma}{original lemmas}
\Crefname{originallemma}{Original lemma}{Original lemmas}
\makeatletter
\newlabel{lemma-etale-makes-quasi-finite-finite-variant}{{00UL}{source}{}{sourcepointer-2}{}}
\newlabel{lemma-etale-makes-quasi-finite-finite-variant@cref}{{[originallemma][0][]00UL}{[1][0][]source}{}{}{}}
\makeatother
\crefname{originallemma}{original lemma}{original lemmas}
\Crefname{originallemma}{Original lemma}{Original lemmas}
\makeatletter
\newlabel{lemma-integral-going-up}{{00GU}{source}{}{sourcepointer-3}{}}
\newlabel{lemma-integral-going-up@cref}{{[originallemma][0][]00GU}{[1][0][]source}{}{}{}}
\makeatother
\crefname{originallemma}{original lemma}{original lemmas}
\Crefname{originallemma}{Original lemma}{Original lemmas}
\makeatletter
\newlabel{lemma-unique-prime-over-localize-below}{{00EA}{source}{}{sourcepointer-4}{}}
\newlabel{lemma-unique-prime-over-localize-below@cref}{{[originallemma][0][]00EA}{[1][0][]source}{}{}{}}
\makeatother
\crefname{originallemma}{original lemma}{original lemmas}
\Crefname{originallemma}{Original lemma}{Original lemmas}
\makeatletter
\newlabel{lemma-unramified-at-prime}{{00UW}{source}{}{sourcepointer-5}{}}
\newlabel{lemma-unramified-at-prime@cref}{{[originallemma][0][]00UW}{[1][0][]source}{}{}{}}
\makeatother
\crefname{originallemma}{original lemma}{original lemmas}
\Crefname{originallemma}{Original lemma}{Original lemmas}
\makeatletter
\newlabel{lemma-unramified-quasi-finite}{{02UR}{source}{}{sourcepointer-6}{}}
\newlabel{lemma-unramified-quasi-finite@cref}{{[originallemma][0][]02UR}{[1][0][]source}{}{}{}}
\makeatother
\title{After an etale neighbourhood the points of an unramified fibre separate into closed pieces}
\author{The Stacks Project Authors}
\date{stacks-project:algebra.tex:local-snapshot-20261012; GNU FDL 1.2 or later}
\begin{document}
\maketitle
\noindent\textit{Editorial note.} All delivered mathematical text is the author's own, copied from the pinned Stacks Project chapter file algebra.tex; no statement or proof step was rewritten, reordered or removed.\par\medskip
\noindent\textit{Editorial note.} Preamble provenance, as required for this source: the frozen snapshot stages the chapter's own preamble file beside it (preamble.tex, 217 lines, 8116 bytes, the same shared Stacks-wide preamble the published chow and etale-cohomology cards use); the first line of the pinned chapter file is \textbackslash{}input\{preamble\}, so no preamble text lives inside the chapter itself. The card takes that file's lines 8-18 (the verbatim package and the three \textbackslash{}comment-based environments reference/slogan/history at lines 9-12, section 13.2, together with the Xy-pic setup) and its lines 154 to the runtime end of the file (the native theorem declarations and mathematical macros). Not one character was changed, no package was added, no macro was defined, and the excerpt serves only as preamble, never as a span coordinate.\par\medskip
\noindent\textit{Editorial note.} The same-chapter pointers are located at their \textbackslash{}label lines inside the pinned chapter file; the cross-chapter pointers are located in the project's own tag list staged in the snapshot (sources/stacks-chow/tags.txt), at the verbatim lines carrying their canonical tags.\par\medskip
\noindent\textit{Editorial note.} The counted statement contains a slogan environment, which the chapter's own preamble defines as a comment environment; the preamble spans of this card start at line 8 and therefore cover it.\par\medskip
\noindent\textit{Editorial note.} The counted proof refers to same-chapter statements by label; they are recorded as navigation stubs pointed at their \textbackslash{}label lines.\par\medskip
\noindent\textit{Allowed standard prerequisites (editorial).}\begin{itemize}
\item The etale splitting of a quasi-finite map. Invoked at the start of the counted proof (line 42244). Reduces the problem to a product decomposition with finite factors, after an etale base change.
\item Surjectivity of a finite map is an open condition on the base. Used at the end of the counted proof (lines 42265-42267). Inverts one element of the base so that the finite factor becomes a surjection.
\end{itemize}\medskip
\ifdefined\hypertarget\hypertarget{sourcepointer-1}{}\fi
\noindent\textit{Source reference.} Original lemma 00DV: sources/stacks-algebra/pinned-original-source.tex, source lines 3833--3833. Navigation pointer for \textbackslash{}ref\{lemma-NAK\} used inside the counted span (line 42262); the argument that consumes it is retained in the span.\par\medskip
\ifdefined\hypertarget\hypertarget{sourcepointer-2}{}\fi
\noindent\textit{Source reference.} Original lemma 00UL: sources/stacks-algebra/pinned-original-source.tex, source lines 40638--40638. Navigation pointer for \textbackslash{}ref\{lemma-etale-makes-quasi-finite-finite-variant\} used inside the counted span (line 42244); the argument that consumes it is retained in the span.\par\medskip
\ifdefined\hypertarget\hypertarget{sourcepointer-3}{}\fi
\noindent\textit{Source reference.} Original lemma 00GU: sources/stacks-algebra/pinned-original-source.tex, source lines 8070--8070. Navigation pointer for \textbackslash{}ref\{lemma-integral-going-up\} used inside the counted span (line 42260); the argument that consumes it is retained in the span.\par\medskip
\ifdefined\hypertarget\hypertarget{sourcepointer-4}{}\fi
\noindent\textit{Source reference.} Original lemma 00EA: sources/stacks-algebra/pinned-original-source.tex, source lines 10014--10014. Navigation pointer for \textbackslash{}ref\{lemma-unique-prime-over-localize-below\} used inside the counted span (line 42258); the argument that consumes it is retained in the span.\par\medskip
\ifdefined\hypertarget\hypertarget{sourcepointer-5}{}\fi
\noindent\textit{Source reference.} Original lemma 00UW: sources/stacks-algebra/pinned-original-source.tex, source lines 41785--41785. Navigation pointer for \textbackslash{}ref\{lemma-unramified-at-prime\} used inside the counted span (line 42254); the argument that consumes it is retained in the span.\par\medskip
\ifdefined\hypertarget\hypertarget{sourcepointer-6}{}\fi
\noindent\textit{Source reference.} Original lemma 02UR: sources/stacks-algebra/pinned-original-source.tex, source lines 41816--41816. Navigation pointer for \textbackslash{}ref\{lemma-unramified-quasi-finite\} used inside the counted span (line 42252); the argument that consumes it is retained in the span.\par\medskip
\begin{lemma}
\label{lemma-etale-makes-unramified-closed}
\begin{slogan}
In an unramified ring map, one can separate the points in a fiber
by passing to an \'etale neighbourhood.
\end{slogan}
Let $R \to S$ be a ring map.
Let $\mathfrak p$ be a prime of $R$.
If $R \to S$ is unramified then there exist
\begin{enumerate}
\item an \'etale ring map $R \to R'$,
\item a prime $\mathfrak p' \subset R'$ lying over $\mathfrak p$.
\item a product decomposition
$$
R' \otimes_R S = A_1 \times \ldots \times A_n \times B
$$
\end{enumerate}
with the following properties
\begin{enumerate}
\item $R' \to A_i$ is surjective,
\item $\mathfrak p'A_i$ is a prime of $A_i$ lying over $\mathfrak p'$, and
\item there is no prime of $B$ lying over $\mathfrak p'$.
\end{enumerate}
\end{lemma}

\begin{proof}
We may apply Lemma \ref{lemma-etale-makes-quasi-finite-finite-variant}.
Thus, after an \'etale base change,
we may assume that $S = A_1 \times \ldots \times A_n \times B$,
that each $R \to A_i$ is finite with exactly one prime $\mathfrak r_i$
lying over $\mathfrak p$ such that
$\kappa(\mathfrak p) \subset \kappa(\mathfrak r_i)$ is purely inseparable,
and that $R \to B$ is not quasi-finite at any prime lying over $\mathfrak p$.
Since $R \to S$ is quasi-finite (see
Lemma \ref{lemma-unramified-quasi-finite})
we see there is no prime of $B$ lying over $\mathfrak p$.
By Lemma \ref{lemma-unramified-at-prime} we see that
$\kappa(\mathfrak r_i)/\kappa(\mathfrak p)$
is separable hence the trivial field extension, and that
$\mathfrak p(A_i)_{\mathfrak r_i}$ is the maximal ideal.
Also, by Lemma \ref{lemma-unique-prime-over-localize-below}
(which applies to $R \to A_i$ because a finite ring map satisfies going up by
Lemma \ref{lemma-integral-going-up})
we have $(A_i)_{\mathfrak r_i} = (A_i)_{\mathfrak p}$.
It follows from Nakayama's Lemma \ref{lemma-NAK}
that the map of local rings
$R_{\mathfrak p} \to (A_i)_{\mathfrak p} = (A_i)_{\mathfrak r_i}$
is surjective. Since $A_i$ is finite over $R$ we see that there
exists a $f \in R$, $f \not \in \mathfrak p$ such that
$R_f \to (A_i)_f$ is surjective. After replacing $R$ by $R_f$ we win.
\end{proof}

\end{document}
